\section*{Chapter \ref{ch:b3:labour-supply-and-demand} --- Labour Supply and Demand in the Data}

\begin{solution}{exo:b3:labour-supply-and-demand:1}
Germany, by 10.1 points (45.2 to 55.3~\%), just ahead of Korea's 9.6 (47.2 to 56.8). % figdata: women_points
Older women work less than younger ones at every date, so a larger share of women
over 65 lowers the average rate even if each age group's rate rises.
\end{solution}

\begin{solution}{exo:b3:labour-supply-and-demand:2}
$0.33 \times 10~\% = 3.3~\%$ more hours.
\end{solution}

\begin{solution}{exo:b3:labour-supply-and-demand:3}
$-(1 - 0.3) \times 1 - 0.3 \times 2 = -0.7 - 0.6 = -1.3$: substitution effect $-0.7$, \omterm{def:b3:labour-supply-and-demand:hicks-marshall}{scale effect} $-0.6$.
\end{solution}

\begin{solution}{exo:b3:labour-supply-and-demand:4}
Middle-wage occupations rich in \omterm{def:b3:labour-supply-and-demand:tasks}{routine tasks} holding their share while
non-routine ones decline; or shares that shift without regard to the routine
content of jobs.
\end{solution}

\begin{solution}{exo:b3:labour-supply-and-demand:5}
The tested solution is \code{ex05_hours.py}, in the chapter's \code{exercises/} folder.
Hours fell by 0.78~\% a year in Germany (from 1960), 0.68 in Korea (from 1963), 0.63 in France,
0.39 in the United Kingdom, 0.29 in Japan (from 1953) and 0.17 in the United States. The
later starting years are the first the source observes: before them its series grows at
a constant rate, which is interpolation, not measurement. % output: ex05_hours
\end{solution}

\begin{solution}{exo:b3:labour-supply-and-demand:6}
The threshold $\sqrt{2F/a}$ falls, so workers with wages between the new and the old
threshold enter. They work $aw$ hours at their wages, below those of the workers already
in, and the jump at the threshold, $\sqrt{2aF}$, is smaller: average hours among workers fall
while total hours rise.
\end{solution}

\begin{solution}{exo:b3:labour-supply-and-demand:7}
$w_R = \sqrt{2 \times 3{,}200/100} = 8$ dalens an hour; hours at it $100 \times 8 = 800$ a year.
\end{solution}

\begin{solution}{exo:b3:labour-supply-and-demand:8}
The tested solution is \code{ex08_ushape.py}, in the chapter's \code{exercises/} folder.
The square term is positive both years, 2.72 (s.e.\ 0.99) in 2000 and 3.97 (0.89) in 2019, and
the minimum moves from about 16{,}500 to about 11{,}000 dollars a person. % output: ex08_ushape
With an $R^2$ near 0.1 the minimum is poorly determined.
\end{solution}

\begin{solution}{exo:b3:labour-supply-and-demand:9}
No change in the employment of older workers when the pension age, an earnings
test or the reward for working longer changes; no bunching of retirements at the
ages where incentives change.
\end{solution}

\begin{solution}{exo:b3:labour-supply-and-demand:10}
The derivation is in the text. The absolute elasticity $(1 - s)\sigma + s\eta$ has derivative $\eta - \sigma$ in $s$:
a larger labour share raises it if and only if $\eta > \sigma$.
\end{solution}

\begin{solution}{exo:b3:labour-supply-and-demand:11}
Take-home pay rises by $4 \times 1.0/1.5 = 2.67~\%$, labour cost falls by $1.33~\%$, and hours rise by
$0.5 \times 2.67 = 1.33~\%$ (equal to $1.0 \times 1.33$ on the demand side).
\end{solution}

\begin{solution}{exo:b3:labour-supply-and-demand:12}
$Y = \sqrt{16/3} = 2.309$: the low and middle groups earn 0.385, the high group 0.770. Middle workers keep the
middle task at capital's price $\rho$ when $\rho$ exceeds what they would earn in the low task,
$\sqrt{2/(3\rho)}/6$, that is when $\rho^3 > 1/54$, or $\rho > 0.265$; above $1/3$ capital is not used at all.
\end{solution}

\begin{solution}{pb:b3:labour-supply-and-demand:hours}
% figdata: brenland_stats
\textbf{1.} 28.1~\% of heads work no hours; among heads under 65, 5.1~\% (94.9~\% work).

\textbf{2.} 1{,}652 hours a year.

\textbf{3.} 16.89 dalens an hour.

\textbf{4.} 0.013 (s.e.\ 0.010): no measurable relation.

\textbf{5.} Workers with higher wages differ in skills and tastes for work; only those
who work are in the sample, so it is selected; and differences in wages between
people are permanent, mixing income and substitution effects, whereas the
elasticity of the chapter is the response of one person to a change in her wage.

\textbf{6.} Hours measured with error appear on both sides: an overstated $h$ raises log
hours and lowers the computed log wage, pulling the slope toward $-1$.

\textbf{7.} $0.33 \times 4 = 1.32~\%$.

\textbf{8.} $0.25 \times 4 = 1.0~\%$: $0.01 \times 4{,}366 = 43.7$ thousand workers.

\textbf{9.} The \omterm{def:b3:labour-supply-and-demand:frisch}{intensive margin}: 1.32~\% more hours from those already working, against
1.0~\% more workers, who enter near the threshold and work fewer hours than the
average.

\textbf{10.} Frictions: adjusting hours takes time and attention, and a small gain is not
worth it at once.

\textbf{11.} Compare hours and participation before and after in reform and other
regions: a difference in differences under parallel trends.

\textbf{12.} $s = 0.479$.

\textbf{13.} $-(1 - 0.479) \times 0.7 - 0.479 \times 1.5 = -1.083$.

\textbf{14.} Substitution $-0.365$, scale $-0.719$.

\textbf{15.} The \omterm{def:b3:labour-supply-and-demand:hicks-marshall}{scale effect}: a firm selling into a competitive world market faces a
very elastic product demand, so $\eta$ is large.

\textbf{16.} Estimates vary mostly with the method and the setting, and published ones
are inflated: use a range, and the rules, not one number.

\textbf{17.} $5 \times 1.083/(0.58 + 1.083) = 3.256~\%$.

\textbf{18.} $5 - 3.256 = 1.744~\%$.

\textbf{19.} Workers: supply (0.58) is less elastic than demand (1.08), and the less elastic
side gets the larger share of a cut in the wedge.

\textbf{20.} $0.58 \times 3.256 = 1.889~\%$.

\textbf{21.} Extensive $0.25 \times 3.256 = 0.814~\%$; intensive $0.33 \times 3.256 = 1.075~\%$.

\textbf{22.} $0.00814 \times 4{,}366 = 35.5$ thousand workers.

\textbf{23.} $0.01075 \times 4{,}366 = 46.9$ thousand full-time equivalents.

\textbf{24.} Workers would move hours into the years of the cut: the \omterm{def:b3:labour-supply-and-demand:frisch}{Frisch elasticity}
applies, and the intensive response is larger while it lasts; firms, which hire for
the long run, may respond less.

\textbf{25.} 35.5 thousand more workers, and 46.9 thousand full-time equivalents of longer
hours among those already working. The two do not add up to a count of full-time
equivalents: question 9 found that entrants work fewer hours than the average, so
82.5 thousand, the sum, is an upper bound, reached only if each entrant worked full time.
The figures also use compensated elasticities for an uncompensated cut, which ignores
the income effect of the higher take-home pay; with leisure a normal good, the
intensive response is smaller.
\end{solution}
